% !TeX root = ../教材题库_全部条目_转写工作稿.tex
% Source: XJ-B1-C02-SU-P060-Y1169
% 第2章 / 小结与复习 小结与复习 / 教材 p.60 / PDF 第67页
% OCR未校对：公式与文字均可能有误；未完成数学TeX转写；禁止直接发给学生
% 原题及相邻原解截图（核对用）：../题目原图/XJ-B1-C02-SU-P060-Y1169-01.webp
% 本文件只在转写完成后作为题干使用；OCR 条目尚不是可靠题干。
\begin{bankitem}{第2章 小结与复习 第3题}{第2章 / 小结与复习 小结与复习 / 教材 p.60 / PDF 第67页}{XJ-B1-C02-SU-P060-Y1169}{OCR未校对：公式与文字均可能有误；未完成数学TeX转写；禁止直接发给学生}
\begin{ocrdraft}
[教材p.60 / PDF第67页]
\par
3．基本不等式以及不等式的性质是证明不等式的重要手段和方法．你能证明
\par
基本不等式吗？在使用基本不等式解决最大（小）值等实际问题时，需要注意什么？
\end{ocrdraft}
\end{bankitem}
